\begin{register}{H}{RMT\_CH\regindex{n}DATA\_REG (\regindex{n} = 0, 1)}{0x0000, 0x0004}\label{regdesc:RMTCHNDATAREG}
\regfield{RMT\_CH\regindex{n}DATA}{32}{0}{{0x000000}}%
\reglabel{Reset}\regnewline%
\begin{regdesc}\begin{reglist}
\label{fielddesc:RMTCHNDATA}\item [RMT\_CH\regindex{n}DATA] 通道 \regindex{n} 通过 APB FIFO 进行读写操作时用到的数据寄存器。(RO)
\end{reglist}\end{regdesc}
\end{register}

\begin{register}{H}{RMT\_CH\regindex{n}DATA\_REG (\regindex{m} = 2, 3)}{0x0008, 0x000C}\label{regdesc:RMTCHMDATAREG}
\regfield{RMT\_CH\regindex{m}DATA}{32}{0}{{0x000000}}%
\reglabel{Reset}\regnewline%
\begin{regdesc}\begin{reglist}
\label{fielddesc:RMTCHMDATA}\item [RMT\_CH\regindex{m}DATA] 通道 \regindex{m} 通过 APB FIFO 进行读写操作时用到的数据寄存器。(RO)
\end{reglist}\end{regdesc}
\end{register}


\begin{register}{H}{RMT\_CH\regindex{n}CONF0\_REG (\regindex{n} = 0, 1)}{0x0010, 0x0014}\label{regdesc:RMTCHNCONF0REG}
\regfield{(reserved)}{7}{25}{0000000}%
\regfield{RMT\_CONF\_UPDATE\_CH\regindex{n}}{1}{24}{{0}}%
\regfield{(reserved)}{1}{23}{0}%
\regfield{RMT\_CARRIER\_OUT\_LV\_CH\regindex{n}}{1}{22}{{1}}%
\regfield{RMT\_CARRIER\_EN\_CH\regindex{n}}{1}{21}{{1}}%
\regfield{RMT\_CARRIER\_EFF\_EN\_CH\regindex{n}}{1}{20}{{1}}%
\regfield{(reserved)}{1}{19}{0}%
\regfield{RMT\_MEM\_SIZE\_CH\regindex{n}}{3}{16}{{0x1}}%
\regfield{RMT\_DIV\_CNT\_CH\regindex{n}}{8}{8}{{0x2}}%
\regfield{RMT\_TX\_STOP\_CH\regindex{n}}{1}{7}{{0}}%
\regfield{RMT\_IDLE\_OUT\_EN\_CH\regindex{n}}{1}{6}{{0}}%
\regfield{RMT\_IDLE\_OUT\_LV\_CH\regindex{n}}{1}{5}{{0}}%
\regfield{RMT\_MEM\_TX\_WRAP\_EN\_CH\regindex{n}}{1}{4}{{0}}%
\regfield{RMT\_TX\_CONTI\_MODE\_CH\regindex{n}}{1}{3}{{0}}%
\regfield{RMT\_APB\_MEM\_RST\_CH\regindex{n}}{1}{2}{{0}}%
\regfield{RMT\_MEM\_RD\_RST\_CH\regindex{n}}{1}{1}{{0}}%
\regfield{RMT\_TX\_START\_CH\regindex{n}}{1}{0}{{0}}%
\reglabel{Reset}\regnewline%
\begin{regdesc}\begin{reglist}
\label{fielddesc:RMTTXSTARTCHN}\item [RMT\_TX\_START\_CH\regindex{n}] 置位此位，通道 \regindex{n} 开始发送数据。(WT)
\label{fielddesc:RMTMEMRDRSTN}\item [RMT\_MEM\_RD\_RST\_CH\regindex{n}] 置位此位，则通道 \regindex{n} 通过发射器访问的 RAM 读地址将被复位。(WT)
\label{fielddesc:RMTAPBMEMRSTCHN}\item [RMT\_APB\_MEM\_RST\_CH\regindex{n}] 置位此位，则通道 \regindex{n} 通过 APB FIFO 访问的读/写 RAM 地址将被复位。(WT)
\label{fielddesc:RMTTXCONTIMODECHN}\item [RMT\_TX\_CONTI\_MODE\_CH\regindex{n}] 
置位此位，使能通道 \regindex{n} 的持续发送模式。在这种模式下，发射器从第一个数据开始发送，如果遇到结束标志，则再次发送第一个数据；如果没有遇到结束标志，发送完最后一个数据后，回卷到第一个数据再次继续发送。(R/W)
\label{fielddesc:RMTMEMTXWRAPENCHN}\item [RMT\_MEM\_TX\_WRAP\_EN\_CH\regindex{n}] 置位此位，使能通道 \regindex{n} 的乒乓发送模式。在这种模式下，如果待发送的数据长度大于该通道的 RAM Block 长度，则发射器将继续从第一个数据开始循环发送。(R/W)
\label{fielddesc:RMTIDLEOUTLVCHN}\item [RMT\_IDLE\_OUT\_LV\_CH\regindex{n}] 配置通道 \regindex{n} 在空闲模式下的输出信号电平。(R/W)
\label{fielddesc:RMTIDLEOUTENCHN}\item [RMT\_IDLE\_OUT\_EN\_CH\regindex{n}] 通道 \regindex{n} 在空闲状态下的输出使能控制位。(R/W)
\label{fielddesc:RMTTXSTOPCHN}\item [RMT\_TX\_STOP\_CH\regindex{n}] 置位此位，则通道 \regindex{n} 的发射器停止发送数据。(R/W/SC)
\label{fielddesc:RMTDIVCNTCHN}\item [RMT\_DIV\_CNT\_CH\regindex{n}] 配置通道 \regindex{n} 的时钟分频器。(R/W)
\label{fielddesc:RMTMEMSIZECHN}\item [RMT\_MEM\_SIZE\_CH\regindex{n}] 配置通道 \regindex{n} 可用的最大 RAM Block 数量。(R/W)
\label{fielddesc:RMTCARRIEREFFENCHN}\item [RMT\_CARRIER\_EFF\_EN\_CH\regindex{n}] 
1：配置通道 \regindex{n} 仅在发送数据状态下对输出信号载波调制；0：配置通道 \regindex{n} 对发送数据状态和空闲状态均加载载波。仅在 RMT\_CARRIER\_EN\_CH\regindex{n} 为 1 时有效。(R/W)
\label{fielddesc:RMTCARRIERENCHN}\item [RMT\_CARRIER\_EN\_CH\regindex{n}] 通道 \regindex{n} 的载波调制使能控制位。1：对输出信号进行载波调制；0：禁止对输出信号进行载波调制。(R/W)
\label{fielddesc:RMTCARRIEROUTLVCHN}\item [RMT\_CARRIER\_OUT\_LV\_CH\regindex{n}] 
配置通道 \regindex{n} 的载波调制方式。(R/W)

1'h0：载波加载在低电平上；

1'h1：载波加载在高电平上。
\label{fielddesc:RMTCONFUPDATECHN}\item [RMT\_CONF\_UPDATE\_CH\regindex{n}] 通道 \regindex{n} 的同步位。 (WT)
\end{reglist}\end{regdesc}
\end{register}


\begin{register}{H}{RMT\_CH\regindex{m}CONF0\_REG (\regindex{m} = 2, 3)}{0x0018, 0x0020}\label{regdesc:RMTCHMCONF0REG}
\regfield{(reserved)}{2}{30}{00}%
\regfield{RMT\_CARRIER\_OUT\_LV\_CH\regindex{m}}{1}{29}{{1}}%
\regfield{RMT\_CARRIER\_EN\_CH\regindex{m}}{1}{28}{{1}}%
\regfield{(reserved)}{2}{26}{00}%
\regfield{RMT\_MEM\_SIZE\_CH\regindex{m}}{3}{23}{{0x1}}%
\regfield{RMT\_IDLE\_THRES\_CH\regindex{m}}{15}{8}{{0x7fff}}%
\regfield{RMT\_DIV\_CNT\_CH\regindex{m}}{8}{0}{{0x2}}%
\reglabel{Reset}\regnewline%
\begin{regdesc}\begin{reglist}
\label{fielddesc:RMTDIVCNTCH2}\item [RMT\_DIV\_CNT\_CH\regindex{m}] 配置通道 \regindex{m} 的时钟分频器。(R/W)
\label{fielddesc:RMTIDLETHRESCH2}\item [RMT\_IDLE\_THRES\_CH\regindex{m}] 配置接收阈值。接收器长时间检测不到信号变化，且持续的时间大于 RMT\_IDLE\_THRES\_CH\regindex{m} 的值，则接收器停止接收过程。(R/W)
\label{fielddesc:RMTMEMSIZECH2}\item [RMT\_MEM\_SIZE\_CH\regindex{m}] 配置通道 \regindex{m} 可用的最大 RAM Block 数量。(R/W)
\label{fielddesc:RMTCARRIERENCH2}\item [RMT\_CARRIER\_EN\_CH\regindex{m}] 通道 \regindex{m} 的载波调制使能控制位。1：对输出信号进行载波调制和解调；0：禁止对输出信号进行载波调制和解调。(R/W)
\label{fielddesc:RMTCARRIEROUTLVCH2}\item [RMT\_CARRIER\_OUT\_LV\_CH\regindex{m}] 配置通道 \regindex{m} 的载波调制方式。(R/W)

1'h0：载波加载在低电平上；

1'h1：载波加载在高电平上。 
\end{reglist}\end{regdesc}
\end{register}


\begin{register}{H}{RMT\_CH\regindex{m}CONF1\_REG (\regindex{m} = 2, 3)}{0x001C, 0x0024}\label{regdesc:RMTCHMCONF1REG}
\regfield{(reserved)}{16}{16}{0000000000000000}%
\regfield{RMT\_CONF\_UPDATE\_CH\regindex{m}}{1}{15}{{0}}%
\regfield{(reserved)}{1}{14}{0}%
\regfield{RMT\_MEM\_RX\_WRAP\_EN\_CH\regindex{m}}{1}{13}{{0}}%
\regfield{RMT\_RX\_FILTER\_THRES\_CH\regindex{m}}{8}{5}{{0xf}}%
\regfield{RMT\_RX\_FILTER\_EN\_CH\regindex{m}}{1}{4}{{0}}%
\regfield{RMT\_MEM\_OWNER\_CH\regindex{m}}{1}{3}{{1}}%
\regfield{RMT\_APB\_MEM\_RST\_CH\regindex{m}}{1}{2}{{0}}%
\regfield{RMT\_MEM\_WR\_RST\_CH\regindex{m}}{1}{1}{{0}}%
\regfield{RMT\_RX\_EN\_CH\regindex{m}}{1}{0}{{0}}%
\reglabel{Reset}\regnewline%
\begin{regdesc}\begin{reglist}
\label{fielddesc:RMTRXENCH2}\item [RMT\_RX\_EN\_CH\regindex{m}] 置位此位，使能通道 \regindex{m} 的接收器，接收器开始接收数据。(R/W)
\label{fielddesc:RMTMEMWRRSTCH2}\item [RMT\_MEM\_WR\_RST\_CH\regindex{m}] 
置位此位，则通道 \regindex{m} 通过接收器访问的 RAM 写地址将被复位。(WT)
\label{fielddesc:RMTAPBMEMRSTCH2}\item [RMT\_APB\_MEM\_RST\_CH\regindex{m}] 置位此位，则通道 \regindex{m} 通过 APB FIFO 访问的读/写 RAM 地址将被复位。(WT)
\label{fielddesc:RMTMEMOWNERCH2}\item [RMT\_MEM\_OWNER\_CH\regindex{m}] 标志通道 \regindex{m} 的 RAM 使用权。(R/W/SC)

1'h1：接收器有权使用该通道的 RAM；

1'h0：APB 总线有权使用该通道的 RAM。
\label{fielddesc:RMTRXFILTERENCH2}\item [RMT\_RX\_FILTER\_EN\_CH\regindex{m}] 通道 \regindex{m} 的接收器滤波功能使能位。(R/W)
\label{fielddesc:RMTRXFILTERTHRESCH2}\item [RMT\_RX\_FILTER\_THRES\_CH\regindex{m}] 接收模式下，忽略宽度小于 RMT\_RX\_FILTER\_THRES\_CH\regindex{m} 个 rmt\_sclk 周期的输入脉冲。(R/W)
\label{fielddesc:RMTMEMRXWRAPENCH2}\item [RMT\_MEM\_RX\_WRAP\_EN\_CH\regindex{m}] 置位此位，使能通道 \regindex{m} 的乒乓接收模式。在这种模式下，如果待接收的数据长度大于该通道的 RAM Block 长度，则接收器将继续把接收数据存入 RAM 的第一个地址，依次循环。(R/W)
\label{fielddesc:RMTCONFUPDATECH2}\item [RMT\_CONF\_UPDATE\_CH\regindex{m}] 通道 \regindex{m} 的同步位。(WT)
\end{reglist}\end{regdesc}
\end{register}



\begin{register}{H}{RMT\_SYS\_CONF\_REG}{0x0068}\label{regdesc:RMTSYSCONFREG}
\regfield{RMT\_CLK\_EN}{1}{31}{{0}}%
\regfield{(reserved)}{4}{27}{0000}%
\regfield{RMT\_SCLK\_ACTIVE}{1}{26}{{1}}%
\regfield{RMT\_SCLK\_SEL}{2}{24}{{0x1}}%
\regfield{RMT\_SCLK\_DIV\_B}{6}{18}{{0x0}}%
\regfield{RMT\_SCLK\_DIV\_A}{6}{12}{{0x0}}%
\regfield{RMT\_SCLK\_DIV\_NUM}{8}{4}{{0x1}}%
\regfield{RMT\_MEM\_FORCE\_PU}{1}{3}{{0}}%
\regfield{RMT\_MEM\_FORCE\_PD}{1}{2}{{0}}%
\regfield{RMT\_MEM\_CLK\_FORCE\_ON}{1}{1}{{0}}%
\regfield{RMT\_APB\_FIFO\_MASK}{1}{0}{{0}}%
\reglabel{Reset}\regnewline%
\begin{regdesc}\begin{reglist}
\label{fielddesc:RMTAPBFIFOMASK}\item [RMT\_APB\_FIFO\_MASK] 1'h1：直接访问 RAM；1'h0：通过 FIFO 访问 RAM。(R/W)
\label{fielddesc:RMTMEMCLKFORCEON}\item [RMT\_MEM\_CLK\_FORCE\_ON] 置位此位，使能 RMT 的 RAM 时钟。(R/W)
\label{fielddesc:RMTMEMFORCEPD}\item [RMT\_MEM\_FORCE\_PD] 置位此位，关闭 RMT RAM。(R/W)
\label{fielddesc:RMTMEMFORCEPU}\item [RMT\_MEM\_FORCE\_PU] 1：禁用 RMT RAM 的 Light-sleep 低功耗模式；0：RMT 处于 Light-sleep 模式时，关闭 RMT RAM。(R/W)
\label{fielddesc:RMTSCLKDIVNUM}\item [RMT\_SCLK\_DIV\_NUM] 小数分频器的整数部分。(R/W)
\label{fielddesc:RMTSCLKDIVA}\item [RMT\_SCLK\_DIV\_A] 小数分频器的分子。(R/W)
\label{fielddesc:RMTSCLKDIVB}\item [RMT\_SCLK\_DIV\_B] 小数分频器的分母。(R/W)
\label{fielddesc:RMTSCLKSEL}\item [RMT\_SCLK\_SEL] 设置 rmt\_sclk 的时钟源：1：APB\_CLK；2：RC\_FAST\_CLK；3：XTAL\_CLK。 (R/W)
\label{fielddesc:RMTSCLKACTIVE}\item [RMT\_SCLK\_ACTIVE] rmt\_sclk 控制开关。(R/W)
\label{fielddesc:RMTCLKEN}\item [RMT\_CLK\_EN] RMT 寄存器的时钟门控使能位。1：打开寄存器的驱动时钟；0：关闭寄存器的驱动时钟。(R/W)
\end{reglist}\end{regdesc}
\end{register}


\begin{register}{H}{RMT\_REF\_CNT\_RST\_REG}{0x0070}\label{regdesc:RMTREFCNTRSTREG}
\regfield{(reserved)}{28}{4}{0000000000000000000000000000}%
\regfield{RMT\_REF\_CNT\_RST\_CH3}{1}{3}{{0}}%
\regfield{RMT\_REF\_CNT\_RST\_CH2}{1}{2}{{0}}%
\regfield{RMT\_REF\_CNT\_RST\_CH1}{1}{1}{{0}}%
\regfield{RMT\_REF\_CNT\_RST\_CH0}{1}{0}{{0}}%
\reglabel{Reset}\regnewline%
\begin{regdesc}\begin{reglist}
\label{fielddesc:RMTREFCNTRSTCH0}\item [RMT\_REF\_CNT\_RST\_CH0] 复位通道 0 的时钟分频器。(WT)
\label{fielddesc:RMTREFCNTRSTCH1}\item [RMT\_REF\_CNT\_RST\_CH1] 复位通道 1 的时钟分频器。(WT)
\label{fielddesc:RMTREFCNTRSTCH2}\item [RMT\_REF\_CNT\_RST\_CH2] 复位通道 2 的时钟分频器。(WT)
\label{fielddesc:RMTREFCNTRSTCH3}\item [RMT\_REF\_CNT\_RST\_CH3] 复位通道 3 的时钟分频器。(WT)
\end{reglist}\end{regdesc}
\end{register}


\begin{register}{H}{RMT\_CH\regindex{n}STATUS\_REG (\regindex{n} = 0, 1)}{0x0028, 0x002C}\label{regdesc:RMTCHNSTATUSREG}
\regfield{RMT\_APB\_MEM\_RADDR\_CH\regindex{n}}{8}{24}{{0x0}}%
\regfield{RMT\_APB\_MEM\_WR\_ERR\_CH\regindex{n}}{1}{23}{{0}}%
\regfield{RMT\_MEM\_EMPTY\_CH\regindex{n}}{1}{22}{{0}}%
\regfield{RMT\_APB\_MEM\_RD\_ERR\_CH\regindex{n}}{1}{21}{{0}}%
\regfield{RMT\_APB\_MEM\_WADDR\_CH\regindex{n}}{9}{12}{{0}}%
\regfield{RMT\_STATE\_CH\regindex{n}}{3}{9}{{0}}%
\regfield{RMT\_MEM\_RADDR\_EX\_CH\regindex{n}}{9}{0}{{0}}%
\reglabel{Reset}\regnewline%
\begin{regdesc}\begin{reglist}
\label{fielddesc:RMTMEMRADDREXCHN}\item [RMT\_MEM\_RADDR\_EX\_CH\regindex{n}] 记录通道 \regindex{n} 发射器使用 RAM 时的地址偏移量。(RO)
\label{fielddesc:RMTSTATECHN}\item [RMT\_STATE\_CH\regindex{n}] 记录通道 \regindex{n} 的 FSM 状态。(RO)
\label{fielddesc:RMTAPBMEMWADDRCHN}\item [RMT\_APB\_MEM\_WADDR\_CH\regindex{n}] 记录 RMT 使用 APB 总线访问 RAM 时的地址偏移量。(RO)
\label{fielddesc:RMTAPBMEMRDERRCHN}\item [RMT\_APB\_MEM\_RD\_ERR\_CH\regindex{n}] RMT 使用 APB 总线进行读 RAM 操作时，如果偏移地址溢出 RAM Block，则该状态位将被置位。(RO)
\label{fielddesc:RMTMEMEMPTYCHN}\item [RMT\_MEM\_EMPTY\_CH\regindex{n}] 发送的数据长度大于 RAM Block 且乒乓模式未启用时，该状态位将被置位。(RO)
\label{fielddesc:RMTAPBMEMWRERRCHN}\item [RMT\_APB\_MEM\_WR\_ERR\_CH\regindex{n}] RMT 使用 APB 总线进行写 RAM 操作时，如果偏移地址溢出 RAM Block，则该状态位将被置位。(RO)
\label{fielddesc:RMTAPBMEMRADDRCHN}\item [RMT\_APB\_MEM\_RADDR\_CH\regindex{n}] 记录 RMT 使用 APB 总线读 RAM 时的地址偏移量。(RO)
\end{reglist}\end{regdesc}
\end{register}

\begin{register}{H}{RMT\_CH\regindex{m}STATUS\_REG (\regindex{m} = 2, 3)}{0x0030, 0x0034}\label{regdesc:RMTCHMSTATUSREG}
\regfield{(reserved)}{4}{28}{0000}%
\regfield{RMT\_APB\_MEM\_RD\_ERR\_CH\regindex{m}}{1}{27}{{0}}%
\regfield{RMT\_MEM\_FULL\_CH\regindex{m}}{1}{26}{{0}}%
\regfield{RMT\_MEM\_OWNER\_ERR\_CH\regindex{m}}{1}{25}{{0}}%
\regfield{RMT\_STATE\_CH\regindex{m}}{3}{22}{{0}}%
\regfield{(reserved)}{1}{21}{0}%
\regfield{RMT\_APB\_MEM\_RADDR\_CH\regindex{m}}{9}{12}{{0}}%
\regfield{(reserved)}{3}{9}{000}%
\regfield{RMT\_MEM\_WADDR\_EX\_CH\regindex{m}}{9}{0}{{0}}%
\reglabel{Reset}\regnewline%
\begin{regdesc}\begin{reglist}
\label{fielddesc:RMTMEMWADDREXCH2}\item [RMT\_MEM\_WADDR\_EX\_CH\regindex{m}] 记录通道 \regindex{m} 接收器使用 RAM 时的地址偏移量。(RO)
\label{fielddesc:RMTAPBMEMRADDRCH2}\item [RMT\_APB\_MEM\_RADDR\_CH\regindex{m}] 记录 RMT 使用 APB 总线访问 RAM 时的地址偏移量。(RO)
\label{fielddesc:RMTSTATECH2}\item [RMT\_STATE\_CH\regindex{m}] 记录通道 \regindex{m} 的 FSM 状态。(RO)
\label{fielddesc:RMTMEMOWNERERRCH2}\item [RMT\_MEM\_OWNER\_ERR\_CH\regindex{m}] RAM Block 使用权发生错误时，该状态位将被置位。(RO)
\label{fielddesc:RMTMEMFULLCH2}\item [RMT\_MEM\_FULL\_CH\regindex{m}] 接收器接收的数据长度大于 RAM Block 时，此状态位将被置位。(RO)
\label{fielddesc:RMTAPBMEMRDERRCH2}\item [RMT\_APB\_MEM\_RD\_ERR\_CH\regindex{m}] RMT 使用 APB 总线执行 RAM 读操作时，如果偏移地址溢出 RAM Block，则该状态位将被置位。(RO)
\end{reglist}\end{regdesc}
\end{register}

\begin{register}{H}{RMT\_INT\_RAW\_REG}{0x0038}\label{regdesc:RMTINTRAWREG}
\regfield{(reserved)}{18}{14}{000000000000000000}%
\regfield{RMT\_CH1\_TX\_LOOP\_INT\_RAW}{1}{13}{{0}}%
\regfield{RMT\_CH0\_TX\_LOOP\_INT\_RAW}{1}{12}{{0}}%
\regfield{RMT\_CH3\_RX\_THR\_EVENT\_INT\_RAW}{1}{11}{{0}}%
\regfield{RMT\_CH2\_RX\_THR\_EVENT\_INT\_RAW}{1}{10}{{0}}%
\regfield{RMT\_CH1\_TX\_THR\_EVENT\_INT\_RAW}{1}{9}{{0}}%
\regfield{RMT\_CH0\_TX\_THR\_EVENT\_INT\_RAW}{1}{8}{{0}}%
\regfield{RMT\_CH3\_ERR\_INT\_RAW}{1}{7}{{0}}%
\regfield{RMT\_CH2\_ERR\_INT\_RAW}{1}{6}{{0}}%
\regfield{RMT\_CH1\_ERR\_INT\_RAW}{1}{5}{{0}}%
\regfield{RMT\_CH0\_ERR\_INT\_RAW}{1}{4}{{0}}%
\regfield{RMT\_CH3\_RX\_END\_INT\_RAW}{1}{3}{{0}}%
\regfield{RMT\_CH2\_RX\_END\_INT\_RAW}{1}{2}{{0}}%
\regfield{RMT\_CH1\_TX\_END\_INT\_RAW}{1}{1}{{0}}%
\regfield{RMT\_CH0\_TX\_END\_INT\_RAW}{1}{0}{{0}}%
\reglabel{Reset}\regnewline%
\begin{regdesc}\begin{reglist}
\label{fielddesc:RMTCH0TXENDINTRAW}\item [RMT\_CH0\_TX\_END\_INT\_RAW] \hyperref[int:RMTCHnTXENDINT]{RMT\_CH0\_TX\_END\_INT} 的原始中断位。(R/WTC/SS)
\label{fielddesc:RMTCH1TXENDINTRAW}\item[RMT\_CH1\_TX\_END\_INT\_RAW] \hyperref[int:RMTCHnTXENDINT]{RMT\_CH1\_TX\_END\_INT} 的原始中断位。(R/WTC/SS)
\label{fielddesc:RMTCH2RXENDINTRAW}\item [RMT\_CH2\_RX\_END\_INT\_RAW] \hyperref[int:RMTCHnRXENDINT]{RMT\_CH2\_RX\_END\_INT} 的原始中断位。(R/WTC/SS)
\label{fielddesc:RMTCH3RXENDINTRAW}\item [RMT\_CH3\_RX\_END\_INT\_RAW] \hyperref[int:RMTCHnRXENDINT]{RMT\_CH3\_RX\_END\_INT} 的原始中断位。(R/WTC/SS)
\label{fielddesc:RMTCH0ERRINTRAW}\item [RMT\_CH0\_ERR\_INT\_RAW] \hyperref[int:RMTCHnERRINT]{RMT\_CH0\_ERR\_INT} 的原始中断位。(R/WTC/SS)
\label{fielddesc:RMTCH1ERRINTRAW}\item [RMT\_CH1\_ERR\_INT\_RAW] \hyperref[int:RMTCHnERRINT]{RMT\_CH1\_ERR\_INT} 的原始中断位。(R/WTC/SS)
\label{fielddesc:RMTCH2ERRINTRAW}\item [RMT\_CH2\_ERR\_INT\_RAW] \hyperref[int:RMTCHnERRINT]{RMT\_CH2\_ERR\_INT} 的原始中断位。(R/WTC/SS)
\label{fielddesc:RMTCH3ERRINTRAW}\item [RMT\_CH3\_ERR\_INT\_RAW] \hyperref[int:RMTCHnERRINT]{RMT\_CH3\_ERR\_INT} 的原始中断位。(R/WTC/SS)
\label{fielddesc:RMTCH0TXTHREVENTINTRAW}\item [RMT\_CH0\_TX\_THR\_EVENT\_INT\_RAW] \hyperref[int:RMTCHnTXTHREVENTINT]{RMT\_CH0\_TX\_THR\_EVENT\_INT} 的原始中断位。(R/WTC/SS)
\label{fielddesc:RMTCH1TXTHREVENTINTRAW}\item [RMT\_CH1\_TX\_THR\_EVENT\_INT\_RAW] \hyperref[int:RMTCHnTXTHREVENTINT]{RMT\_CH1\_TX\_THR\_EVENT\_INT} 的原始中断位。(R/WTC/SS)
\label{fielddesc:RMTCH2RXTHREVENTINTRAW}\item [RMT\_CH2\_RX\_THR\_EVENT\_INT\_RAW] \hyperref[int:RMTCHmRXTHREVENTINT]{RMT\_CH2\_RX\_THR\_EVENT\_INT} 的原始中断位。(R/WTC/SS)
\label{fielddesc:RMTCH3RXTHREVENTINTRAW}\item [RMT\_CH3\_RX\_THR\_EVENT\_INT\_RAW] \hyperref[int:RMTCHmRXTHREVENTINT]{RMT\_CH3\_RX\_THR\_EVENT\_INT} 的原始中断位。(R/WTC/SS)
\label{fielddesc:RMTCH0TXLOOPINTRAW}\item [RMT\_CH0\_TX\_LOOP\_INT\_RAW] \hyperref[int:RMTCHnTXLOOPINT]{RMT\_CH0\_TX\_LOOP\_INT} 的原始中断位。(R/WTC/SS)
\label{fielddesc:RMTCH1TXLOOPINTRAW}\item [RMT\_CH1\_TX\_LOOP\_INT\_RAW] \hyperref[int:RMTCHnTXLOOPINT]{RMT\_CH1\_TX\_LOOP\_INT} 的原始中断位。(R/WTC/SS)
\end{reglist}\end{regdesc}
\end{register}


\begin{register}{H}{RMT\_INT\_ST\_REG}{0x003C}\label{regdesc:RMTINTSTREG}
\regfield{(reserved)}{18}{14}{000000000000000000}%
\regfield{RMT\_CH1\_TX\_LOOP\_INT\_ST}{1}{13}{{0}}%
\regfield{RMT\_CH0\_TX\_LOOP\_INT\_ST}{1}{12}{{0}}%
\regfield{RMT\_CH3\_RX\_THR\_EVENT\_INT\_ST}{1}{11}{{0}}%
\regfield{RMT\_CH2\_RX\_THR\_EVENT\_INT\_ST}{1}{10}{{0}}%
\regfield{RMT\_CH1\_TX\_THR\_EVENT\_INT\_ST}{1}{9}{{0}}%
\regfield{RMT\_CH0\_TX\_THR\_EVENT\_INT\_ST}{1}{8}{{0}}%
\regfield{RMT\_CH3\_ERR\_INT\_ST}{1}{7}{{0}}%
\regfield{RMT\_CH2\_ERR\_INT\_ST}{1}{6}{{0}}%
\regfield{RMT\_CH1\_ERR\_INT\_ST}{1}{5}{{0}}%
\regfield{RMT\_CH0\_ERR\_INT\_ST}{1}{4}{{0}}%
\regfield{RMT\_CH3\_RX\_END\_INT\_ST}{1}{3}{{0}}%
\regfield{RMT\_CH2\_RX\_END\_INT\_ST}{1}{2}{{0}}%
\regfield{RMT\_CH1\_TX\_END\_INT\_ST}{1}{1}{{0}}%
\regfield{RMT\_CH0\_TX\_END\_INT\_ST}{1}{0}{{0}}%
\reglabel{Reset}\regnewline%
\begin{regdesc}\begin{reglist}
\label{fielddesc:RMTCH0TXENDINTST}\item [RMT\_CH0\_TX\_END\_INT\_ST] \hyperref[int:RMTCHnTXENDINT]{RMT\_CH0\_TX\_END\_INT} 的屏蔽中断状态位。(RO)
\label{fielddesc:RMTCH1TXENDINTST}\item [RMT\_CH1\_TX\_END\_INT\_ST] \hyperref[int:RMTCHnTXENDINT]{RMT\_CH1\_TX\_END\_INT} 的屏蔽中断状态位。(RO)
\label{fielddesc:RMTCH2RXENDINTST}\item [RMT\_CH2\_RX\_END\_INT\_ST] \hyperref[int:RMTCHnRXENDINT]{RMT\_CH2\_RX\_END\_INT} 的屏蔽中断状态位。(RO)
\label{fielddesc:RMTCH3RXENDINTST}\item [RMT\_CH3\_RX\_END\_INT\_ST] \hyperref[int:RMTCHnRXENDINT]{RMT\_CH3\_RX\_END\_INT} 的屏蔽中断状态位。(RO)
\label{fielddesc:RMTCH0ERRINTST}\item [RMT\_CH0\_ERR\_INT\_ST] \hyperref[int:RMTCHnERRINT]{RMT\_CH0\_ERR\_INT} 的屏蔽中断状态位。(RO)
\label{fielddesc:RMTCH1ERRINTST}\item [RMT\_CH1\_ERR\_INT\_ST] \hyperref[int:RMTCHnERRINT]{RMT\_CH1\_ERR\_INT} 的屏蔽中断状态位。(RO)
\label{fielddesc:RMTCH2ERRINTST}\item [RMT\_CH2\_ERR\_INT\_ST] \hyperref[int:RMTCHnERRINT]{RMT\_CH2\_ERR\_INT} 的屏蔽中断状态位。(RO)
\label{fielddesc:RMTCH3ERRINTST}\item [RMT\_CH3\_ERR\_INT\_ST] \hyperref[int:RMTCHnERRINT]{RMT\_CH3\_ERR\_INT} 的屏蔽中断状态位。(RO)
\label{fielddesc:RMTCH0TXTHREVENTINTST}\item [RMT\_CH0\_TX\_THR\_EVENT\_INT\_ST] \hyperref[int:RMTCHnTXTHREVENTINT]{RMT\_CH0\_TX\_THR\_EVENT\_INT} 的屏蔽中断状态位。(RO)
\label{fielddesc:RMTCH1TXTHREVENTINTST}\item [RMT\_CH1\_TX\_THR\_EVENT\_INT\_ST] \hyperref[int:RMTCHnTXTHREVENTINT]{RMT\_CH1\_TX\_THR\_EVENT\_INT} 的屏蔽中断状态位。(RO)
\label{fielddesc:RMTCH2RXTHREVENTINTST}\item [RMT\_CH2\_RX\_THR\_EVENT\_INT\_ST] \hyperref[int:RMTCHmRXTHREVENTINT]{RMT\_CH2\_RX\_THR\_EVENT\_INT} 的屏蔽中断状态位。(RO)
\label{fielddesc:RMTCH3RXTHREVENTINTST}\item [RMT\_CH3\_RX\_THR\_EVENT\_INT\_ST] \hyperref[int:RMTCHmRXTHREVENTINT]{RMT\_CH3\_RX\_THR\_EVENT\_INT} 的屏蔽中断状态位。(RO)
\label{fielddesc:RMTCH0TXLOOPINTST}\item [RMT\_CH0\_TX\_LOOP\_INT\_ST] \hyperref[int:RMTCHnTXLOOPINT]{RMT\_CH0\_TX\_LOOP\_INT} 的屏蔽中断状态位。(RO)
\label{fielddesc:RMTCH1TXLOOPINTST}\item [RMT\_CH1\_TX\_LOOP\_INT\_ST] \hyperref[int:RMTCHnTXLOOPINT]{RMT\_CH1\_TX\_LOOP\_INT} 的屏蔽中断状态位。(RO)
\end{reglist}\end{regdesc}
\end{register}


\begin{register}{H}{RMT\_INT\_ENA\_REG}{0x0040}\label{regdesc:RMTINTENAREG}
\regfield{(reserved)}{18}{14}{000000000000000000}%
\regfield{RMT\_CH1\_TX\_LOOP\_INT\_ENA}{1}{13}{{0}}%
\regfield{RMT\_CH0\_TX\_LOOP\_INT\_ENA}{1}{12}{{0}}%
\regfield{RMT\_CH3\_RX\_THR\_EVENT\_INT\_ENA}{1}{11}{{0}}%
\regfield{RMT\_CH2\_RX\_THR\_EVENT\_INT\_ENA}{1}{10}{{0}}%
\regfield{RMT\_CH1\_TX\_THR\_EVENT\_INT\_ENA}{1}{9}{{0}}%
\regfield{RMT\_CH0\_TX\_THR\_EVENT\_INT\_ENA}{1}{8}{{0}}%
\regfield{RMT\_CH3\_ERR\_INT\_ENA}{1}{7}{{0}}%
\regfield{RMT\_CH2\_ERR\_INT\_ENA}{1}{6}{{0}}%
\regfield{RMT\_CH1\_ERR\_INT\_ENA}{1}{5}{{0}}%
\regfield{RMT\_CH0\_ERR\_INT\_ENA}{1}{4}{{0}}%
\regfield{RMT\_CH3\_RX\_END\_INT\_ENA}{1}{3}{{0}}%
\regfield{RMT\_CH2\_RX\_END\_INT\_ENA}{1}{2}{{0}}%
\regfield{RMT\_CH1\_TX\_END\_INT\_ENA}{1}{1}{{0}}%
\regfield{RMT\_CH0\_TX\_END\_INT\_ENA}{1}{0}{{0}}%
\reglabel{Reset}\regnewline%
\begin{regdesc}\begin{reglist}
\label{fielddesc:RMTCH0TXENDINTENA}\item [RMT\_CH0\_TX\_END\_INT\_ENA] \hyperref[int:RMTCHnTXENDINT]{RMT\_CH0\_TX\_END\_INT} 的中断使能位。(R/W)
\label{fielddesc:RMTCH1TXENDINTENA}\item [RMT\_CH1\_TX\_END\_INT\_ENA] \hyperref[int:RMTCHnTXENDINT]{RMT\_CH1\_TX\_END\_INT} 的中断使能位。(R/W)
\label{fielddesc:RMTCH2RXENDINTENA}\item [RMT\_CH2\_RX\_END\_INT\_ENA] \hyperref[int:RMTCHnRXENDINT]{RMT\_CH2\_RX\_END\_INT} 的中断使能位。(R/W)
\label{fielddesc:RMTCH3RXENDINTENA}\item [RMT\_CH3\_RX\_END\_INT\_ENA] \hyperref[int:RMTCHnRXENDINT]{RMT\_CH3\_RX\_END\_INT} 的中断使能位。(R/W)
\label{fielddesc:RMTCH0ERRINTENA}\item [RMT\_CH0\_ERR\_INT\_ENA] \hyperref[int:RMTCHnERRINT]{RMT\_CH0\_ERR\_INT} 的中断使能位。(R/W)
\label{fielddesc:RMTCH1ERRINTENA}\item [RMT\_CH1\_ERR\_INT\_ENA] \hyperref[int:RMTCHnERRINT]{RMT\_CH1\_ERR\_INT} 的中断使能位。(R/W)
\label{fielddesc:RMTCH2ERRINTENA}\item [RMT\_CH2\_ERR\_INT\_ENA] \hyperref[int:RMTCHnERRINT]{RMT\_CH2\_ERR\_INT} 的中断使能位。(R/W)
\label{fielddesc:RMTCH3ERRINTENA}\item [RMT\_CH3\_ERR\_INT\_ENA] \hyperref[int:RMTCHnERRINT]{RMT\_CH3\_ERR\_INT} 的中断使能位。(R/W)
\label{fielddesc:RMTCH0TXTHREVENTINTENA}\item [RMT\_CH0\_TX\_THR\_EVENT\_INT\_ENA] \hyperref[int:RMTCHnTXTHREVENTINT]{RMT\_CH0\_TX\_THR\_EVENT\_INT} 的中断使能位。(R/W)
\label{fielddesc:RMTCH1TXTHREVENTINTENA}\item [RMT\_CH1\_TX\_THR\_EVENT\_INT\_ENA] \hyperref[int:RMTCHnTXTHREVENTINT]{RMT\_CH1\_TX\_THR\_EVENT\_INT} 的中断使能位。(R/W)
\label{fielddesc:RMTCH2RXTHREVENTINTENA}\item [RMT\_CH2\_RX\_THR\_EVENT\_INT\_ENA] \hyperref[int:RMTCHmRXTHREVENTINT]{RMT\_CH2\_RX\_THR\_EVENT\_INT} 的中断使能位。(R/W)
\label{fielddesc:RMTCH3RXTHREVENTINTENA}\item [RMT\_CH3\_RX\_THR\_EVENT\_INT\_ENA] \hyperref[int:RMTCHmRXTHREVENTINT]{RMT\_CH3\_RX\_THR\_EVENT\_INT} 的中断使能位。(R/W)
\label{fielddesc:RMTCH0TXLOOPINTENA}\item [RMT\_CH0\_TX\_LOOP\_INT\_ENA] \hyperref[int:RMTCHnTXLOOPINT]{RMT\_CH0\_TX\_LOOP\_INT} 的中断使能位。(R/W)
\label{fielddesc:RMTCH1TXLOOPINTENA}\item [RMT\_CH1\_TX\_LOOP\_INT\_ENA] \hyperref[int:RMTCHnTXLOOPINT]{RMT\_CH1\_TX\_LOOP\_INT} 的中断使能位。(R/W)
\end{reglist}\end{regdesc}
\end{register}


\begin{register}{H}{RMT\_INT\_CLR\_REG}{0x0044}\label{regdesc:RMTINTCLRREG}
\regfield{(reserved)}{18}{14}{000000000000000000}%
\regfield{RMT\_CH1\_TX\_LOOP\_INT\_CLR}{1}{13}{{0}}%
\regfield{RMT\_CH0\_TX\_LOOP\_INT\_CLR}{1}{12}{{0}}%
\regfield{RMT\_CH3\_RX\_THR\_EVENT\_INT\_CLR}{1}{11}{{0}}%
\regfield{RMT\_CH2\_RX\_THR\_EVENT\_INT\_CLR}{1}{10}{{0}}%
\regfield{RMT\_CH1\_TX\_THR\_EVENT\_INT\_CLR}{1}{9}{{0}}%
\regfield{RMT\_CH0\_TX\_THR\_EVENT\_INT\_CLR}{1}{8}{{0}}%
\regfield{RMT\_CH3\_ERR\_INT\_CLR}{1}{7}{{0}}%
\regfield{RMT\_CH2\_ERR\_INT\_CLR}{1}{6}{{0}}%
\regfield{RMT\_CH1\_ERR\_INT\_CLR}{1}{5}{{0}}%
\regfield{RMT\_CH0\_ERR\_INT\_CLR}{1}{4}{{0}}%
\regfield{RMT\_CH3\_RX\_END\_INT\_CLR}{1}{3}{{0}}%
\regfield{RMT\_CH2\_RX\_END\_INT\_CLR}{1}{2}{{0}}%
\regfield{RMT\_CH1\_TX\_END\_INT\_CLR}{1}{1}{{0}}%
\regfield{RMT\_CH0\_TX\_END\_INT\_CLR}{1}{0}{{0}}%
\reglabel{Reset}\regnewline%
\begin{regdesc}\begin{reglist}
\label{fielddesc:RMTCH0TXENDINTCLR}\item [RMT\_CH0\_TX\_END\_INT\_CLR] \hyperref[int:RMTCHnTXENDINT]{RMT\_CH0\_TX\_END\_INT} 的中断清除位。(WT)
\label{fielddesc:RMTCH1TXENDINTCLR}\item [RMT\_CH1\_TX\_END\_INT\_CLR] \hyperref[int:RMTCHnTXENDINT]{RMT\_CH1\_TX\_END\_INT} 的中断清除位。(WT)
\label{fielddesc:RMTCH2RXENDINTCLR}\item [RMT\_CH2\_RX\_END\_INT\_CLR] \hyperref[int:RMTCHnRXENDINT]{RMT\_CH2\_RX\_END\_INT} 的中断清除位。(WT)
\label{fielddesc:RMTCH3RXENDINTCLR}\item [RMT\_CH3\_RX\_END\_INT\_CLR] \hyperref[int:RMTCHnRXENDINT]{RMT\_CH3\_RX\_END\_INT} 的中断清除位。(WT)
\label{fielddesc:RMTCH0ERRINTCLR}\item [RMT\_CH0\_ERR\_INT\_CLR] \hyperref[int:RMTCHnERRINT]{RMT\_CH0\_ERR\_INT} 的中断清除位。(WT)
\label{fielddesc:RMTCH1ERRINTCLR}\item [RMT\_CH1\_ERR\_INT\_CLR] \hyperref[int:RMTCHnERRINT]{RMT\_CH1\_ERR\_INT} 的中断清除位。(WT)
\label{fielddesc:RMTCH2ERRINTCLR}\item [RMT\_CH2\_ERR\_INT\_CLR] \hyperref[int:RMTCHnERRINT]{RMT\_CH2\_ERR\_INT} 的中断清除位。(WT)
\label{fielddesc:RMTCH3ERRINTCLR}\item [RMT\_CH3\_ERR\_INT\_CLR] \hyperref[int:RMTCHnERRINT]{RMT\_CH3\_ERR\_INT} 的中断清除位。(WT)
\label{fielddesc:RMTCH0TXTHREVENTINTCLR}\item [RMT\_CH0\_TX\_THR\_EVENT\_INT\_CLR] \hyperref[int:RMTCHnTXTHREVENTINT]{RMT\_CH0\_TX\_THR\_EVENT\_INT} 的中断清除位。(WT)
\label{fielddesc:RMTCH1TXTHREVENTINTCLR}\item [RMT\_CH1\_TX\_THR\_EVENT\_INT\_CLR] \hyperref[int:RMTCHnTXTHREVENTINT]{RMT\_CH1\_TX\_THR\_EVENT\_INT} 的中断清除位。(WT)
\label{fielddesc:RMTCH2RXTHREVENTINTCLR}\item [RMT\_CH2\_RX\_THR\_EVENT\_INT\_CLR] \hyperref[int:RMTCHmRXTHREVENTINT]{RMT\_CH2\_RX\_THR\_EVENT\_INT} 的中断清除位。(WT)
\label{fielddesc:RMTCH3RXTHREVENTINTCLR}\item [RMT\_CH3\_RX\_THR\_EVENT\_INT\_CLR] \hyperref[int:RMTCHmRXTHREVENTINT]{RMT\_CH3\_RX\_THR\_EVENT\_INT} 的中断清除位。(WT)
\label{fielddesc:RMTCH0TXLOOPINTCLR}\item [RMT\_CH0\_TX\_LOOP\_INT\_CLR] \hyperref[int:RMTCHnTXLOOPINT]{RMT\_CH0\_TX\_LOOP\_INT} 的中断清除位。(WT)
\label{fielddesc:RMTCH1TXLOOPINTCLR}\item [RMT\_CH1\_TX\_LOOP\_INT\_CLR] \hyperref[int:RMTCHnTXLOOPINT]{RMT\_CH1\_TX\_LOOP\_INT} 的中断清除位。(WT)
\end{reglist}\end{regdesc}
\end{register}


\begin{register}{H}{RMT\_CH\regindex{n}CARRIER\_DUTY\_REG (\regindex{n} = 0, 1)}{0x0048, 0x004C}\label{regdesc:RMTCHNCARRIERDUTYREG}
\regfield{RMT\_CARRIER\_HIGH\_CH\regindex{n}}{16}{16}{{0x40}}%
\regfield{RMT\_CARRIER\_LOW\_CH\regindex{n}}{16}{0}{{0x40}}%
\reglabel{Reset}\regnewline%
\begin{regdesc}\begin{reglist}
\label{fielddesc:RMTCARRIERLOWCHN}\item [RMT\_CARRIER\_LOW\_CH\regindex{n}] 配置通道 \regindex{n} 载波的低电平时钟周期。(R/W)
\label{fielddesc:RMTCARRIERHIGHCHN}\item [RMT\_CARRIER\_HIGH\_CH\regindex{n}] 配置通道 \regindex{n} 载波的高电平时钟周期。(R/W)
\end{reglist}\end{regdesc}
\end{register}


\begin{register}{H}{RMT\_CH\regindex{m}\_RX\_CARRIER\_RM\_REG (\regindex{m} = 2, 3)}{0x0050, 0x0054}\label{regdesc:RMTCHMRXCARRIERRMREG}
\regfield{RMT\_CARRIER\_HIGH\_THRES\_CH\regindex{m}}{16}{16}{{0x00}}%
\regfield{RMT\_CARRIER\_LOW\_THRES\_CH\regindex{m}}{16}{0}{{0x00}}%
\reglabel{Reset}\regnewline%
\begin{regdesc}\begin{reglist}
\label{fielddesc:RMTCARRIERLOWTHRESCH2}\item [RMT\_CARRIER\_LOW\_THRES\_CH\regindex{m}] 载波调制模式下，通道 \regindex{m} 低电平周期为 RMT\_CARRIER\_LOW\_THRES\_CH\regindex{m} + 1。(R/W)
\label{fielddesc:RMTCARRIERHIGHTHRESCH2}\item [RMT\_CARRIER\_HIGH\_THRES\_CH\regindex{m}] 载波调制模式下，通道 \regindex{m} 高电平周期为 RMT\_CARRIER\_HIGH\_THRES\_CH\regindex{m} + 1。(R/W)
\end{reglist}\end{regdesc}
\end{register}

\begin{register}{H}{RMT\_CH\regindex{n}\_TX\_LIM\_REG (\regindex{n} = 0, 1)}{0x0058, 0x005C}\label{regdesc:RMTCHNTXLIMREG}
\regfield{(reserved)}{11}{21}{00000000000}%
\regfield{RMT\_LOOP\_COUNT\_RESET\_CH\regindex{n}}{1}{20}{{0}}%
\regfield{RMT\_TX\_LOOP\_CNT\_EN\_CH\regindex{n}}{1}{19}{{0}}%
\regfield{RMT\_TX\_LOOP\_NUM\_CH\regindex{n}}{10}{9}{{0}}%
\regfield{RMT\_TX\_LIM\_CH\regindex{n}}{9}{0}{{0x80}}%
\reglabel{Reset}\regnewline%
\begin{regdesc}\begin{reglist}
\label{fielddesc:RMTTXLIMCHN}\item [RMT\_TX\_LIM\_CH\regindex{n}] 配置通道 \regindex{n} 发送脉冲编码数量的上限值。(R/W)
\label{fielddesc:RMTTXLOOPNUMCHN}\item [RMT\_TX\_LOOP\_NUM\_CH\regindex{n}] 配置持续发送模式下最大循环发送次数。(R/W)
\label{fielddesc:RMTTXLOOPCNTENCHN}\item [RMT\_TX\_LOOP\_CNT\_EN\_CH\regindex{n}] 置位此位，使能循环次数计数。(R/W)
\label{fielddesc:RMTLOOPCOUNTRESETCHN}\item [RMT\_LOOP\_COUNT\_RESET\_CH\regindex{n}] 重置持续发送模式下的循环计数器。(WT)
\end{reglist}\end{regdesc}
\end{register}


\begin{register}{H}{RMT\_TX\_SIM\_REG}{0x006C}\label{regdesc:RMTTXSIMREG}
\regfield{(reserved)}{29}{3}{00000000000000000000000000000}%
\regfield{RMT\_TX\_SIM\_EN}{1}{2}{{0}}%
\regfield{RMT\_TX\_SIM\_CH1}{1}{1}{{0}}%
\regfield{RMT\_TX\_SIM\_CH0}{1}{0}{{0}}%
\reglabel{Reset}\regnewline%
\begin{regdesc}\begin{reglist}
\label{fielddesc:RMTTXSIMCH0}\item [RMT\_TX\_SIM\_CH0] 置位此位，使能通道 0 与其它启用的通道同步开始发送数据。 (R/W)
\label{fielddesc:RMTTXSIMCH1}\item [RMT\_TX\_SIM\_CH1] 置位此位，使能通道 1 与其它启用的通道同步开始发送数据。 (R/W)
\label{fielddesc:RMTTXSIMEN}\item [RMT\_TX\_SIM\_EN] 置位此位，多个通道开始同步发送数据。(R/W)
\end{reglist}\end{regdesc}
\end{register}


\begin{register}{H}{RMT\_CH\regindex{m}\_RX\_LIM\_REG (\regindex{m} = 2, 3)}{0x0060, 0x0064}\label{regdesc:RMTCHMRXLIMREG}
\regfield{(reserved)}{23}{9}{00000000000000000000000}%
\regfield{RMT\_CH\regindex{m}\_RX\_LIM\_REG}{9}{0}{{0x80}}%
\reglabel{Reset}\regnewline%
\begin{regdesc}\begin{reglist}
\label{fielddesc:RMTRXLIMCH2}\item [RMT\_RX\_LIM\_CH\regindex{m}] 配置通道 \regindex{m} 最大可接收的脉冲编码数量。(R/W)
\end{reglist}\end{regdesc}
\end{register}


\begin{register}{H}{RMT\_DATE\_REG}{0x00CC}\label{regdesc:RMTDATEREG}
\regfield{(reserved)}{4}{28}{0000}%
\regfield{RMT\_RMT\_DATE}{28}{0}{{0x2006231}}%
\reglabel{Reset}\regnewline%
\begin{regdesc}\begin{reglist}
\label{fielddesc:RMTDATE}\item [RMT\_DATE] 版本控制寄存器。(R/W)
\end{reglist}\end{regdesc}
\end{register}
